\section{Conceptos básicos: Clasificación de los sistemas Agentic}

\subsection{¿Qué es un Agent?}

El término «Agent» tiene múltiples definiciones en la industria:
\begin{itemize}
    \item \textbf{Sentido amplio}: Sistema completamente autónomo capaz de ejecutarse de forma independiente durante periodos prolongados y utilizar diversas herramientas para completar tareas complejas.
    \item \textbf{Sentido estricto}: Implementación programática que se ejecuta siguiendo un flujo de trabajo predefinido.
\end{itemize}

Anthropic agrupa todas estas variantes bajo el nombre de \textbf{Agentic Systems} (sistemas Agentic), pero establece una distinción clave a nivel de arquitectura:

\begin{importantbox}{Distinción central entre Workflow y Agent}
\begin{itemize}
    \item \textbf{Workflow} (flujo de trabajo): Orquestación mediante \textbf{rutas de código predefinidas} para el LLM y las herramientas. El flujo de control lo deciden los desarrolladores.
    \item \textbf{Agent} (agente autónomo): El LLM decide \textbf{de forma dinámica y autónoma} su flujo de procesamiento y el uso de herramientas. Se entrega el control al propio modelo.
\end{itemize}
\end{importantbox}

La importancia de esta distinción radica en que los Workflows son más predecibles y controlables, lo que los hace adecuados para tareas estructuradas; mientras que los Agentes son más flexibles, ideales para problemas abiertos. Ambos no son opuestos, sino posiciones diferentes dentro del espectro de complejidad.

\subsection{Cuándo utilizar (y cuándo no) los sistemas Agentic}

Anthropic proporciona un marco de decisión claro:

\begin{knowledgebox}{Ruta de decisión con incremento progresivo de la complejidad}
\begin{enumerate}
    \item \textbf{Primero considere la solución más simple}: ¿Son suficientes una única llamada al LLM + Recuperación (Retrieval) + Ejemplos en contexto (In-context Examples)?
    \item \textbf{Se requiere mayor fiabilidad}: Utilice el modo Workflow para obtener predictibilidad y consistencia.
    \item \textbf{Se necesita flexibilidad y toma de decisiones impulsada por el modelo}: Utilice el modo Agent.
\end{enumerate}
\end{knowledgebox}

\begin{warningbox}{Trampa de introducir complejidad prematuramente}
Los sistemas Agentic suelen obtener un mejor desempeño en las tareas a cambio de \textbf{mayor latencia y costo}. Introducir arquitecturas Agentes complejas desde las primeras etapas del proyecto es un error común; en muchos escenarios, una única llamada al LLM cuidadosamente optimizada puede ser totalmente suficiente. Es fundamental evaluar primero el límite superior de las soluciones simples antes de decidir adoptar modos más complejos.
\end{warningbox}

\subsection{Recomendaciones sobre el uso de marcos de trabajo}

Existen en el mercado muchos marcos que simplifican el desarrollo de sistemas agentic, como el Claude Agent SDK, el AWS Strands Agents SDK, Rivet (constructor de flujos de trabajo con interfaz gráfica por arrastrar y soltar), entre otros. La recomendación de Anthropic respecto a los marcos es:

\begin{itemize}
    \item \textbf{Comience con las APIs nativas}: Muchos patrones pueden implementarse con apenas unas pocas líneas de código.
    \item \textbf{Comprensione la implementación subyacente}: La capa de abstracción del marco puede ocultar detalles del prompt y la respuesta, aumentando la dificultad para depurar (debug).
    \item \textbf{Evite la complejidad innecesaria}: Los marcos pueden inducir a los desarrolladores a añadir complejidad innecesaria cuando una solución simple es suficiente.
\end{itemize}

\begin{warningbox}{Errores comunes en el uso de marcos}
Las \textbf{suposiciones erróneas} sobre la implementación subyacente son una causa frecuente de errores por parte del cliente. Al utilizar un marco, asegúrese de comprender exactamente qué hace dicho marco a nivel inferior: ¿qué prompt envía? ¿Cómo parsea las llamadas a herramientas (tool calls)? ¿Cómo se transmiten los resultados de los pasos intermedios? Sin entender esto, no podrá depurar eficazmente.
\end{warningbox}

\subsection{Resumen de la sección}
Los sistemas Agentic se dividen en dos grandes categorías: Workflow y Agent. El flujo del Workflow lo controla el código, mientras que el flujo del Agente lo controla el modelo. Al desarrollar, debe partir de soluciones simples e incrementar la complejidad únicamente cuando sea necesario. Los marcos pueden ayudar a iniciar rápidamente, pero es esencial comprender la implementación subyacente.



